\providecommand{\rank}{\operatorname{r}}
\providecommand{\RR}{\mathbb R}
\providecommand{\CC}{\mathbb C}
\providecommand{\KK}{\mathbb K}

% ---------- Exercise 66 ----------
\setcounter{enumi}{65}
\item {\bfseries\boldmath
State, with justification, whether in the real vector space $\RR^3$
}

\begin{enumerate}[label=\alph*)]

\item $(-8,5,-7)$ is a linear combination of $\bigl((2,-1,1),(2,0,-2)\bigr)$.

\begin{answer}
Yes:
$(-8,5,-7)=-5(2,-1,1)+(2,0,-2)$.
\end{answer}

\item $(-8,5,-7)$ is a linear combination of $\bigl((2,-1,1),(1,0,-2)\bigr)$.

\begin{answer}
No. In a putative expression
$(-8,5,-7)=\alpha(2,-1,1)+\beta(1,0,-2)$, the second coordinate gives
$\alpha=-5$; the first then gives $\beta=2$. The third would be
$\alpha-2\beta=-9\ne-7$.
\end{answer}

\item $ (2,-10,-17)$ is a linear combination of $\bigl((4,-6,1),(1,-2,-1),(-1,0,-4)\bigr)$.

\begin{answer}
Yes:
$ (2,-10,-17)=-3(4,-6,1)+14(1,-2,-1)+0(-1,0,-4)$.
\end{answer}

\item $(-5,3,-5)$ is a linear combination of $\bigl((4,-6,1),(1,-2,-1),(-1,0,-4)\bigr)$.

\begin{answer}
No. Each of the three given generators satisfies
$-8x-5y+2z=0$, so every linear combination does too.
For $(-5,3,-5)$ the left side is $40-15-10=15\ne0$.
\end{answer}

\item $(-5,3,-5)$ is a linear combination of $\bigl((1,0,1),(0,1,-1),(1,0,-1),(2,-1,1)\bigr)$.

\begin{answer}
Write $u_1=(1,0,1)$, $u_2=(0,1,-1)$, $u_3=(1,0,-1)$,
$u_4=(2,-1,1)$. Then
\[
 (-5,3,-5)=-\tfrac72u_1+3u_2-\tfrac32u_3+0u_4.
\]
Thus the answer is yes.
\end{answer}

\item $(0,1,-1)$ is a linear combination of $\bigl((1,0,1),(0,1,-1),(1,0,-1),(2,-1,1)\bigr)$.

\begin{answer}
Yes: $(0,1,-1)=0u_1+1u_2+0u_3+0u_4$.
\end{answer}

\item $(1,1,0)$ is a linear combination of $\bigl((1,0,1),(0,1,-1),(1,0,-1),(2,-1,1)\bigr)$.

\begin{answer}
Yes: $(1,1,0)=u_1+u_2$.
\end{answer}

\item $(0,0,0)$ is a linear combination of $\bigl((1,0,1),(0,1,-1),(1,0,-1),(2,-1,1)\bigr)$.

\begin{answer}
Yes: $(0,0,0)=0u_1+0u_2+0u_3+0u_4$.
\end{answer}

\item $(-6,3,12)$ is a linear combination of $(2,-1,-4)$.

\begin{answer}
Yes: $(-6,3,12)=-3(2,-1,-4)$.
\end{answer}

\item $(4,6,-8)$ is a linear combination of $(2,-1,-4)$.

\begin{answer}
No. The first coordinate of $(4,6,-8)=\alpha(2,-1,-4)$ would give
$\alpha=2$, but the second would then be $-2\ne6$.
\end{answer}

\item $\left(\frac{\sqrt5}{\sqrt3},\frac{\sqrt5}{\sqrt2},\frac5{\sqrt2}\right)$ is a linear combination of $\frac15(\sqrt{10},\sqrt{15},5\sqrt3)$.

Is $\frac15(\sqrt{10},\sqrt{15},5\sqrt3)$ a linear combination of $\left(\frac{\sqrt5}{\sqrt3},\frac{\sqrt5}{\sqrt2},\frac5{\sqrt2}\right)$?

\begin{answer}
Put
\[
 u=\left(\frac{\sqrt5}{\sqrt3},\frac{\sqrt5}{\sqrt2},\frac5{\sqrt2}\right),
 \qquad v=\frac15(\sqrt{10},\sqrt{15},5\sqrt3).
\]
Coordinatewise simplification gives
$u=\frac5{\sqrt6}v$ and $v=\frac{\sqrt6}{5}u$.
Both questions therefore have an affirmative answer.
\end{answer}
\end{enumerate}
